\subsection{An ordinary infinite-support ghost in the full class-three image}


\begin{statement}{A.59}{Proved}{A recursive ordinary ghost with an actual class-three lift}
There is an infinite-support ordinary completed unit-equation pair in the
locus $\mathcal V_3$ of A.58. In particular the actual
class-three section map has a section whose abelian
coordinates are ordinary point coordinates at every
finite place and which is neither a rational-point
image nor a cuspidal image.

The pair is constructed by three sequences of distinct
primes $q_j,r_j,s_j>3$. Put
\[
 n_j=(j+4)!,\qquad D_j=48n_j^2,\qquad
 E_j(p)=24n_j(p^2-1).
 \tag{A.59.1}
\]
Its two completed multiplicative coordinates are
\[
 A=3\prod_{j\ge0}\frac{q_j^{E_j(q_j)}}{s_j^{E_j(s_j)}},
 \qquad
 B=-2\prod_{j\ge0}\frac{r_j^{E_j(r_j)}}{s_j^{E_j(s_j)}}.
 \tag{A.59.2}
\]
At every finite prime $p$ they are the completed
classes of ordinary nonzero numbers $x_p,1-x_p\in\Q_p$.
Both completed classes have infinite valuation support.
\end{statement}
\begin{proof}
We prove the prime recursion, including its simultaneous
local constraints and all compatibility with cyclotomic
fields. The class-three assertion then follows by applying
the actual obstruction calculation A.58 to its limit.

\emph{The three radical fields.}
At every stage the coefficients will have the form
\[
 a=3u^{24},\qquad b=-2v^{24},\qquad c=w^{24},
 \tag{A.59.3}
\]
where $u,v,w$ are positive integers supported on three
disjoint sets of primes different from two and three.
In particular $a,b,c$ are pairwise coprime.
For $24\mid n$ put $F_n=\Q(\mu_n)$ and consider
\[
 L_q=F_n((b/c)^{1/n}),\quad
 L_r=F_n((a/c)^{1/n}),\quad
 L_s=F_n((-a/b)^{1/n}).
 \tag{A.59.4}
\]
Each has maximal abelian subfield exactly $F_n$.
Here are the field calculations, including the denominator
in the third radicand. Their squareclasses are respectively
$-2,3,6$, whose square roots belong to $\Q(\mu_{24})$.
Their valuations at two or three include an exponent
equal to $1$ or $-1$. None is an odd-prime power in the
abelian field $F_n$: for an odd prime $\ell\mid n$,
the radicand itself has valuation $1$ or $-1$ at
two or three, so its $\ell$th-root polynomial is
irreducible. Its real root would generate a degree-$\ell$
real field whose other conjugates include nonreal roots.
It is therefore not Galois over $\Q$. Such a field
cannot be an intermediate field of an abelian extension.
No removal of a twenty-fourth-power factor is used
in this odd-prime argument.

Nor is any radicand a fourth power in $F_n$.
Its fourth root, after multiplication by elements
already in $F_n$, would give, respectively,
$2^{1/4},3^{1/4},6^{1/4}$. The exact comparisons are
\[
 \begin{split}
 (b/c)^{1/4}&=(v/w)^6\zeta_8\,2^{1/4},\\
 (a/c)^{1/4}&=(u/w)^6\,3^{1/4},\\
 (-a/b)^{1/4}&=(u/v)^6\,6^{1/4}/\sqrt2 .
 \end{split}
 \tag{A.59.5}
\]
The positive fourth-root fields are real and are not
Galois over $\Q$; the same prime valuations prove
irreducibility. This gives a contradiction to membership
in $F_n$.

Kummer theory now gives a cyclic kernel $N$ of order
$n/2$ in each extension (A.59.4). Explicitly, the order
of a radicand in $F_n^\times/(F_n^\times)^n$ cannot
be smaller than $n/2$: a smaller order would make it
a fourth power or an odd-prime power. The square root
already in $F_n$ gives the opposite bound.
Conjugation by a lift of
$t\in(\Z/n)^\times$ acts on $N$ by multiplication
by $t$. Since $24\mid n$,
\[
       \gcd\bigl(n/2,\ t-1:t\in(\Z/n)^\times\bigr)=2.
 \tag{A.59.6}
\]
All $t$ are odd; choosing $t=-1$ on the two-primary
factor gives precisely one factor of two, and the
Chinese remainder theorem supplies $t\not\equiv1$
at each odd prime divisor. Thus $[N,G]=N^2$.
The fixed field of $N^2$ contains the splitting field
$\Q(d^{1/4},i)$ with $d=2,3,6$, respectively, by
(A.59.5). This is a nonabelian dihedral extension
of order eight. Hence $G/N^2$ is nonabelian,
whereas $G/N$ is abelian; since $N/N^2$ has order
two, it follows that $[G,G]=N$.
This proves the maximal-abelian assertion.

Consequently, for any modulus $M$ divisible by $n$,
each field in (A.59.4) intersects $\Q(\mu_M)$ in
exactly $F_n$. A unit class modulo $M$ which is one
modulo $n$ can therefore be imposed simultaneously
with complete splitting in any one of these three
fields. The relevant automorphism of the compositum
is the identity on the radical field and that unit
class on $\Q(\mu_M)$. Chebotarev supplies infinitely
many primes with those specifications
\cite[Theorem 8.31]{AgentMilneANT}.

\paragraph{Exact local image at the coefficient primes.}
Let $p\mid n$, put $s=v_p(n)$, and assume $s\ge3$ when
$p=2$.  Set
\[
 z_0=1+p^s,\qquad
 \mathcal B=z_0+p^{s+1}\Z_p,\qquad
 F_n(z)=\exp\bigl(24n(z^2-1)\log z\bigr)
                   \quad(z\in1+p^s\Z_p).
\tag{A.59.C1}
\]
This definition fixes the Teichmuller component to one.
For positive rational integers $z$ in this ball it equals
the ordinary integer power $z^{24n(z^2-1)}$.
Put $h=3s+v_p(48)$ and
\[
 \epsilon_p=\begin{cases}1,&p\ne2,\\3,&p=2.\end{cases}
\]
Then the exact image is
\[
                 F_n(\mathcal B)
                   =F_n(z_0)+p^{h+\epsilon_p}\Z_p.
\tag{A.59.C2}
\]
For any $a,b,c\in\Z_p$ of which at least one is a unit,
the image of the three-variable map on $\mathcal B^3$ is
\[
 \{aF_n(q)+bF_n(r)-cF_n(t):q,r,t\in\mathcal B\}
       =(a+b-c)F_n(z_0)+p^{h+\epsilon_p}\Z_p.
\tag{A.59.C3}
\]
Consequently this expression can be made zero, with all
three variables retaining exact valuation
$v_p(q-1)=v_p(r-1)=v_p(t-1)=s$, exactly when
\[
                         v_p(a+b-c)\ge h+\epsilon_p.
\tag{A.59.C4}
\]
One may leave two variables equal to $z_0$ and vary only
a variable with unit coefficient.

\emph{Proof of the local image formula.}
We work in the ring $R=\Z_p\langle w\rangle$ of restricted
power series, whose coefficients tend to zero $p$-adically.
Write $z=1+p^s w$ and set
\[
 A=48np^{2s},\qquad v_p(A)=h.
\]
Expansion of the logarithm gives the exact identity
\[
 \begin{split}
 24n\bigl((1+p^s w)^2-1\bigr)\log(1+p^s w)
      &=A w^2V(w),\\
 V(w)&=1+\sum_{k\ge2}
       \frac{(-1)^k(k-1)p^{sk}}{2k(k+1)}w^k.
 \end{split}
\tag{A.59.C5}
\]
There is no linear term in $V$.  To check the displayed
coefficient, multiply
$\log(1+x)/x=\sum_{k\ge0}(-1)^kx^k/(k+1)$ by
$1+x/2$.  The coefficient of $x^k$ for $k\ge1$ is
$(-1)^k(k-1)/(2k(k+1))$.

For odd $p$, the coefficients of $V-1$ belong to $pR$.
Indeed, for $k\ge2$,
$v_p(k(k+1))\le k-1$, so their valuations are at least
$sk-(k-1)\ge1$.  They tend to infinity since the
denominator valuations grow only logarithmically.
At two the same estimate gives
\[
      v_2\left(\frac{(k-1)2^{sk}}{2k(k+1)}\right)
          \ge k(s-1)\ge2s-2\ge4.
\]
Thus
\[
 V\equiv1\pmod{pR}\quad(p\ne2),\qquad
 V\equiv1\pmod{16R}\quad(p=2).
\tag{A.59.C6}
\]
The assumption $s\ge3$ at two also guarantees that all
logarithms and exponentials in (A.59.C5) converge on this
entire ball.

The exponential series in the complete ring $R$ now
has the form
\[
 F_n(1+p^s w)=1+A f(w),\qquad
 f(w)=w^2V(w)+
       \sum_{j\ge2}\frac{A^{j-1}}{j!}(w^2V(w))^j.
\tag{A.59.C7}
\]
This is a restricted integral series.  For example
$v_p(j!)\le j-1$ shows both integrality and convergence
of the second sum, since $h\ge3$ for odd $p$ and
$h\ge13$ at two.  That sum is in $pR$ for odd $p$
and in $16R$ at two.  Hence
\[
 f(w)=w^2+pR_1(w)\quad(p\ne2),\qquad
 f(w)=w^2+16R_1(w)\quad(p=2)
\tag{A.59.C8}
\]
for an integral restricted series $R_1$.

An integral restricted series is one-Lipschitz on
$\Z_p$, and its difference under $w\mapsto1+p^r t$
is divisible by $p^r$ in $\Z_p\langle t\rangle$.
For odd $p$, (A.59.C8) therefore gives
\[
       \frac{f(1+pt)-f(1)}p=2t+pR_2(t),
                   \qquad R_2\in\Z_p\langle t\rangle.
\tag{A.59.C9}
\]
The right side is a bijection of $\Z_p$: for a specified
value $y$, the equation
$t=(y-pR_2(t))/2$ has a unique solution by contraction,
because two is a unit and $R_2$ is one-Lipschitz.
It follows that
$f(1+p\Z_p)=f(1)+p\Z_p$.

At two, every odd square is one modulo eight, so
(A.59.C8) first gives
\[
                 f(1+2\Z_2)\subset f(1)+8\Z_2.
\]
On the smaller ball $w=1+4t$, the same equation gives
\[
        \frac{f(1+4t)-f(1)}8=t+2R_3(t),
                    \qquad R_3\in\Z_2\langle t\rangle.
\tag{A.59.C10}
\]
To see the divisibility of the error directly, the square
term contributes $t+2t^2$, while
$16(R_1(1+4t)-R_1(1))/8$ is divisible by eight.
The right side of (A.59.C10) is again a bijection of
$\Z_2$, by the contraction equation
$t=y-2R_3(t)$.  Thus the smaller ball already maps onto
$f(1)+8\Z_2$, proving equality for $1+2\Z_2$ as well.

Since $z\in\mathcal B$ corresponds exactly to
$w\in1+p\Z_p$, multiplying these exact images by $A$
in (A.59.C7) proves (A.59.C2).  All changes in the
three-variable expression consequently lie in the
ideal on the right of (A.59.C3).  A variable with unit
coefficient fills that ideal by itself while the
other two variables stay at $z_0$.
This proves (A.59.C3), and (A.59.C4) follows because
$F_n(z_0)$ is a unit.  Membership in $\mathcal B$
retains the required exact valuation $s$ for each
variable.

\paragraph{Precision and the derivative check.}
The derivative with respect to $z$ at $z_0$ is
\[
 \frac{F_n'(z_0)}{F_n(z_0)}
   =24n\left(2z_0\log z_0+\frac{z_0^2-1}{z_0}\right),
 \qquad v_p(F_n'(z_0))=2s+v_p(96).
\tag{A.59.C11}
\]
Indeed the bracket is
$4p^s+(2/3)p^{3s}-(5/6)p^{4s}+\cdots$;
the remainder has strictly larger valuation than
the leading term under the assumptions above.
Equivalently (A.59.C7)--(A.59.C8) give the same valuation
after multiplying the derivative in $w$ by $p^{-s}$.
For the invariant proposed in the recursion,
$H_j(p)=2(2s+v_p(96))+j+3$,
one has $H_j(p)\ge h+\epsilon_p$ at every coefficient
prime.  Thus that invariant suffices for (A.59.C4);
the stronger bound using the derivative in $w$
also suffices.

Finally, approximate any chosen local solution by rational
integers modulo $p^L$, with $L\ge s+1$.  The variables
remain in $\mathcal B$.  Their $w$-coordinates change by
elements of $p^{L-s}\Z_p$; the one-Lipschitz property
of $f$ in (A.59.C7) makes each $F_n$ value change by an
element of $p^{h+L-s}\Z_p$.
Taking $L$ arbitrarily large therefore imposes any
specified next-stage error precision.  At the integer
prime variables supplied by the later CRT and
Chebotarev step, (A.59.C1) is the original integer power,
so this precision statement applies to the actual
prime recursion.


\emph{The all-place invariant.}
Put $B(D)=(D+1)^4$. Define
\[
 d_j(p)=
 \begin{cases}
  2v_p(n_j)+v_p(96),&p\mid n_j,\\
  0,&p\nmid n_j,
 \end{cases}
 \qquad H_j(p)=2d_j(p)+j+3.
 \tag{A.59.7}
\]
We construct coefficients $a_j,b_j,c_j$ of the form
(A.59.3), with pairwise disjoint supports as there,
such that at every finite prime $p$ the homogeneous
equation
\[
       a_jX^{D_j}+b_jY^{D_j}=c_jZ^{D_j}
           \quad\hbox{has }X,Y,Z\in\Z_p^\times.
 \tag{A.59.8}
\]
We also require
\[
 v_p(a_j+b_j-c_j)\ge H_j(p)
                       \qquad(p\le B(D_j)).
 \tag{A.59.9}
\]
The initial values $(a_0,b_0,c_0)=(3,-2,1)$
satisfy both assertions with $(X,Y,Z)=(1,1,1)$.
Maintaining $D_j$-solubility at all places, including
every coefficient-support prime, will remove any
need to keep a separate age classification of
the primes in those supports.

\emph{The triangular congruence construction.}
At stage $j$ write $n=n_j$, $D=D_j$, and
$(a,b,c)=(a_j,b_j,c_j)$. Choose an integer
\[
 P\ge B(D_{j+1})
       \quad\hbox{which is at least every prime in }
                      \operatorname{supp}(abc).
 \tag{A.59.10}
\]
We process every prime $p\le P$ in increasing order.
For each such $p$ choose an integer $L_p$ at least
$H_{j+1}(p)$, with $p^{L_p}>P$; at $p\mid n$
require also $L_p\ge v_p(n)+2$.
These finite choices are made before the residues
are chosen. In particular every prime power
dividing $h-1$ for an integer $h\le P$ is covered
by the prescribed precision at its prime divisor.

For each of the three eventual prime variables,
the residues at primes $p\le P$ will satisfy
\[
 \begin{array}{ll}
 v_p(z-1)=v_p(n),&p\mid n,\\
 z\not\equiv0,1,-1\pmod p,&p\nmid n.
 \end{array}
 \tag{A.59.11}
\]
The second line is imposed only at $p\ge5$.
Before reaching $p$, all prime factors of $p-1$
have already been processed. The preceding
precisions therefore fix, separately for the
three variables, their residues modulo $p-1$
and the exponent classes
\[
             e_i(p)=24n(z_i^2-1)\pmod{p-1},
                         \qquad i=a,b,c.
 \tag{A.59.12}
\]
There is no dependence on an unchosen larger
prime in this definition.

For every prime $\ell\mid p-1$, (A.59.11) gives
the exact valuation
\[
 v_\ell(24n(z_i^2-1))=v_\ell(48n^2).
 \tag{A.59.13}
\]
For odd $\ell\mid n$, the two factors $z_i-1$
and $z_i+1$ have valuations $v_\ell(n)$ and zero.
At two they have valuations $v_2(n)$ and one.
For $\ell\nmid n$, the second line of (A.59.11)
gives valuation zero on both sides of (A.59.13).
It follows that all three power maps on
$\mathbf F_p^\times$ have exactly the same image:
\[
 \gcd(e_i(p),p-1)=\gcd(D,p-1),\qquad
       (\mathbf F_p^\times)^{e_i(p)}
                       =(\mathbf F_p^\times)^D.
 \tag{A.59.14}
\]

We now produce local input values for the three
variables at $p$, and then retain their residues
modulo $p^{L_p}$.
At $p\mid n$ use the balls
$1+p^{v_p(n)}+p^{v_p(n)+1}\Z_p$ from (A.59.C1).
These primes are at most $B(D)$, so (A.59.9)
applies. The estimate in (A.59.C11) gives
$H_j(p)\ge3v_p(n)+v_p(48)+\epsilon_p$.
Since at least one of $a,b,c$ is a unit,
(A.59.C2)--(A.59.C4) supply input values in those
three balls with
\[
                  aF_n(z_a)+bF_n(z_b)-cF_n(z_c)=0.
 \tag{A.59.15}
\]
They retain the first line of (A.59.11).

For $p\nmid n$ we must include the already fixed
Teichmuller exponent (A.59.12). On a residue ball
with $z\equiv u\ne0,1,-1\pmod p$, let $\omega$
be the Teichmuller lift of $u$ and define
\[
 {\cal F}_{i,p}(z)=
     \omega^{e_i(p)}
       \exp\!\left(24n(z^2-1)\log(z/\omega)\right).
 \tag{A.59.16}
\]
For an integer satisfying all the earlier
congruences this is exactly its ordinary power
$z^{24n(z^2-1)}$.
The map (A.59.16) is a bijection from
$\omega(1+p\Z_p)$ onto
$\omega^{e_i(p)}(1+p\Z_p)$.
Here is a direct inverse calculation. Substituting
$z=\omega(1+pt)$ and dividing the logarithm by $p$
gives
\[
 \frac{\log({\cal F}_{i,p}(z)/\omega^{e_i(p)})}{p}
       =24n(\omega^2-1)t+pG(t),
                 \qquad G\in\Z_p\langle t\rangle.
 \tag{A.59.17}
\]
The linear coefficient is a unit. For $k\ge2$,
the coefficient $p^{k-1}/k$ from
$\log(1+pt)/p$ is divisible by $p$.
The integral restricted series $G$ is
one-Lipschitz, so the equation for a specified
right side is a contraction equation for $t$.
This proves bijectivity and, at the same time,
the required local continuity with arbitrary
later precision.

Suppose first that $p\le B(D)$.
Choose one common value
$\lambda\in(\mathbf F_p^\times)^D$
and preimages under each of the three power
maps in (A.59.14), all different from $1,-1$.
If this image group has more than one element,
choose $\lambda\ne1$; an even exponent sends
both $1$ and $-1$ to one. If its image is
$\{1\}$, every nonzero residue is a preimage
and $p\ge5$ leaves a choice different from
$1,-1$. Let $\widetilde\lambda$ be the
Teichmuller lift of $\lambda$.
The three local maps each attain
$\widetilde\lambda(1+p\Z_p)$.
Their common values $\widetilde\lambda$ give
sum-error
$(a+b-c)\widetilde\lambda\in p\Z_p$
by (A.59.9). Varying a variable with unit
coefficient cancels that error within the
same image ball. This gives an exact local
version of (A.59.15).

Suppose next that $p>B(D)$.
Use the unit point (A.59.8), reduce it modulo
$p$, and scale its three coordinates by one
nonzero common residue. This scaling can be
chosen so that none of the three $D$th-power
values is one. Indeed, the group
$(\mathbf F_p^\times)^D$ has cardinality
$(p-1)/\gcd(D,p-1)>3$, and at most three
of its elements are forbidden by the three
coordinates. Lift that scaling to a $p$-adic
unit. The resulting three $D$th-power values
are ordinary $p$-adic units satisfying the
homogeneous equation exactly.
By (A.59.14) choose residue preimages under
the three exponent maps; they avoid $1,-1$.
The bijections (A.59.16)--(A.59.17) lift these
preimages to the three prescribed $p$-adic
values. This proves the required exact local
equation even at a large prime dividing a
coefficient.

At every processed prime we now approximate
these three local input values modulo
$p^{L_p}$. Equations (A.59.C7) and (A.59.17)
show that the resulting power values retain
the error bound
\[
                    v_p(aF(z_a)+bF(z_b)-cF(z_c))
                                      \ge H_{j+1}(p).
 \tag{A.59.18}
\]
The restrictions (A.59.11) are unchanged.
The exponent classes at all subsequent primes
are then defined by these chosen residues,
as in (A.59.12). This completes the triangular
construction.

\emph{The actual primes and the next all-place invariant.}
Put $M=\prod_{p\le P}p^{L_p}$.
Each of the three CRT classes is a unit class
which is one modulo $n$.
Apply (A.59.4)--(A.59.6) and Chebotarev to
choose distinct primes $q,r,s>P$ in these
respective classes, splitting completely in
$L_q,L_r,L_s$, respectively.
All three fields use the old coefficients.
The primes may be chosen successively increasing,
so that they are distinct from one another
and every earlier coefficient prime.
Set
\[
 \begin{split}
 a'&=a q^{24n(q^2-1)},\\
 b'&=b r^{24n(r^2-1)},\\
 c'&=c s^{24n(s^2-1)}.
 \end{split}
 \tag{A.59.19}
\]
Their supports and twenty-fourth-power forms
remain as required. At all $p\le P$,
(A.59.18) gives the next error bound.
It also gives unit $D_{j+1}$-points:
at $(1,1,1)$ vary a variable with unit
coefficient and apply one-variable Hensel,
using
\[
              H_{j+1}(p)>2v_p(D_{j+1}).
 \tag{A.59.20}
\]
This inequality follows directly from (A.59.7).

At each new coefficient prime $Q\in\{q,r,s\}$
one has
\[
                  \gcd(D_{j+1},Q-1)\mid n.
 \tag{A.59.21}
\]
For a prime $\ell\mid n$, its valuation in
$Q-1$ is exactly $v_\ell(n)$.
A prime dividing $D_{j+1}$ but not $n$
is at most $j+5\le P$; (A.59.11) excludes
$Q\equiv1\pmod\ell$.
This proves (A.59.21), and since $n\mid D_{j+1}$,
the $n$th-power and $D_{j+1}$th-power images
in $\mathbf F_Q^\times$ are equal.

At $q$, complete splitting in $L_q$ makes
$c/b$ an $n$th power modulo $q$.
The later factors $r^{24n(r^2-1)}$ and
$s^{24n(s^2-1)}$ are already $n$th powers.
Thus $c'/b'$ is a $D_{j+1}$th power modulo $q$.
Set $X=Z=1$ and solve
$Y^{D_{j+1}}=(c'-a')/b'$ by simple-root
Hensel. This is a unit root because its
residue is $c'/b'$ and $q\nmid D_{j+1}$.
At $r$ set $Y=Z=1$ and solve for $X$ using
the split ratio $c/a$.
At $s$ set $Y=Z=1$ and solve for $X$ using
the split ratio $-b/a$, the inverse of the
third radicand in (A.59.4).
This proves the new unit-point condition at
all three new primes. No mutual congruence
between the new primes was required.

Every remaining prime exceeds $P\ge B(D_{j+1})$,
and all three coefficients are units there.
The nonzero-coordinate finite-field estimate
proved by Jacobi sums in A.44 gives a point
on
$(a'/c')X^{D_{j+1}}+(b'/c')Y^{D_{j+1}}=1$.
Its coordinate derivatives are units, so it
lifts with $X,Y\in\Z_p^\times$ and $Z=1$.
This covers every finite prime.
The next precision requirement (A.59.9)
holds because all primes at most
$B(D_{j+1})$ were processed.
The induction is complete.

\emph{Ordinary limits and the class-three image.}
Name the primes selected at stage $j$ by
$q_j,r_j,s_j$. Every new exponent in (A.59.19)
is divisible by $n_j$, and the factorials
$n_j$ are eventually divisible by any fixed
positive integer. Hence the three products
defining completed classes $a_\infty,b_\infty,c_\infty$,
and the quotients (A.59.2), converge in the
completed multiplicative group.

For a fixed finite prime $p$, each valuation
of $a_j,b_j,c_j$ is eventually constant:
no selected prime is used twice. All later
factors are $p$-adic units whose exponents
are divisible by the growing factorials.
These exponents annihilate every finite
quotient of $\Z_p^\times$ eventually.
Thus the three coefficient sequences
converge to ordinary nonzero elements
$a_p,b_p,c_p\in\Q_p^\times$, and their
completed Kummer classes are the localized
classes $a_\infty,b_\infty,c_\infty$.
For every fixed $p$, the bound (A.59.9)
eventually applies and tends to infinity.
It gives $a_p+b_p=c_p$.
Set
\[
                    x_p=a_p/c_p,\qquad
                  1-x_p=b_p/c_p.
 \tag{A.59.22}
\]
Both entries are ordinary and nonzero.
At the real place the signs of $(A,B)$
are $(+,-)$, represented for example by
the ordinary pair $(3,-2)$.

At two and three the valuations of $(A,B)$
are exactly (A.58.1). At the selected primes
they are
\[
 \begin{array}{c|c}
 q_j&(E_j(q_j),0)\\
 r_j&(0,E_j(r_j))\\
 s_j&(-E_j(s_j),-E_j(s_j)).
 \end{array}
 \tag{A.59.23}
\]
All other valuation pairs are zero.
In the notation of A.58, $e_{q_j}=e_{r_j}=e_{s_j}=n_j$.
This family converges to zero in $\widehat\Z$.
Therefore $(A,B)\in\mathcal V_3$, and A.58
supplies an actual continuous class-three
section above its distinguished class-two
section.

Neither abelian coordinate has finite
valuation support, so it cannot be the
Kummer class of a nonzero rational number.
This excludes rational-point section images.
Cusp zero has trivial second coordinate,
excluded by $v_2(B)=1$; cusp one has
trivial first coordinate, excluded by
$v_3(A)=1$; cusp infinity has equal
valuation coordinates, excluded already
at two. These statements use the explicit
abelian cusp maps and therefore apply to
their images in the class-three quotient.
\end{proof}

The retained certificate
\path{code/triangular_crt_certificate.py}
checks the exact valuations in (A.59.13), the
power-image equality (A.59.14), the availability
of residue preimages different from $1,-1$,
and the fresh-prime divisibility (A.59.21).
Its output records the finite cases.
The independent modular-image certificate
\path{code/class3_coefficient_image_certificate.py}
and its output check the exact local image at
the coefficient primes, including the two-adic
depth in (A.59.C2).
The prime recursion itself is an existence
construction using the stated Chebotarev
applications; no complete large prime-search
stage is claimed as executed.
